\chapter{Sales platform \textsuperscript{Neumayer}}

Parts of the diploma project were the implementation of a web site to sell the application and a way to handle licenses to use the application. The sales platform is split into two parts:

\begin{description}
\item [The frontend:] provides the web site to purchase the license from and a RESTful API that can be accessed using HTTP requests. \hfill \\
\item [The backend:] generates and stores user licenses in a local SQLite database and validates them. \hfill \\
\end{description}

The backend and frontend run on a Linux-based server and communicate through an IPC mechanism called D-Bus.

\section{D-Bus}

Note: The following description is based on information from the FreeDesktop wiki.\cite{dbus1}
\\
D-Bus, short for "Desktop Bus", is an inter-process-communication mechanism used to allow 
processes to call functions of other processes in a secure, predictable, and 
standardized manner. Although it is mostly used in applications with graphical user interfaces 
and desktop environments on Unix-like platforms, D-Bus does not have a functional dependency 
on graphical environments. Implementations of D-Bus are available for Linux, BSD and Windows, 
among other operating systems.
\\
Some of the most important features of D-Bus related to this project which lead to the conclusion that choosing D-Bus as the only viable IPC mechanism, are:
\begin{itemize}
\item {Communication between processes without the need of the network stack of the operating system kernel: the backend doesn't have to listen on a network-based socket. Having two processes on the same machine exchanging information over the network stack of the operating system would unavoidably add overhead to the exchanged messages.}
\item {Influenced by object-oriented programming languages: Methods which are exposed via D-Bus are identifiable by a globally unique name space, which is akin to what object-oriented programming languages provide. This provides a familiar development environment and eases the process of finding errors.}
\item {The ability to introspect exposed methods: it is possible to browse through all exposed methods on the message bus at run time, checking their method names, parameters and return types. This makes the efforts of debugging the communication less complicated.}
\item {Access control via policy handling: Access to the exposed methods can optionally be granted to specific applications only, for example by which user account started the client process. This feature adds an extra layer of security and ensures that information does not leak into unprivileged processes.}
\item {On-demand activation: The server process can be started automatically by the D-Bus daemon as soon as a client process tries to access the exposed methods. This can be accomplished by an XML-based configuration file which makes the executable path of the server process and its exposed methods known to the bus daemon.}
\end{itemize}

\section{Frontend}
To serve web pages to potential customers and provide a REST-ful API, a frontend server runs on the TCP port 80. This server is implemented in the Python programming language and it uses the programming library Bottle for its templating system. Using this templating system it is possible to stitch parts of the website together at runtime for each seperate HTTP request.

\section{Backend}
The backend is a background daemon program which is implemented in \cpp. To store data about customers and their licenses the daemon writes to a SQLite database. Licenses are generated for uses through asymmetrical encryption using APIs provided by OpenSSL.

\section{Separation of privileges}
The frontend and the backend are separated from each other to provide a more secure architecture. As the public way to access the data on the database is through the frontend (via HTTP requests), it is critical to avoid vulnerabilities which could be triggered by third-parties. 
\\
Process-level separation already provides some security as the sales platform is split into two parts and operating systems usually have the required mechanisms built into them to not allow one process to read memory of another process. To further secure the architecture, both processes run as separate users and each user only has access to its own contained directory.

\section{Communication between frontend and backend}
To allow the frontend to return license data to the user it needs a well-defined interface to get access to the data. As mentioned before, the backend provides an interface which can be accessed through the D-Bus IPC mechanism. Only the user of the frontend process has access to this interface as provided by the specific D-Bus policy file. The D-Bus policy file looks like this:

\begin{lstlisting}
<!DOCTYPE busconfig PUBLIC
 "-//freedesktop//DTD D-BUS Bus Configuration 1.0//EN"
 "http://www.freedesktop.org/standards/dbus/1.0/busconfig.dtd">
 
<!-- This file has to be copied to /etc/dbus-1/system.d/ -->
<busconfig>
        <policy user="btsbackend">
                <allow own="at.beatthescore.license"/>
        </policy>
        <policy user="btsfrontend">
		        <allow send_destination="at.beatthescore.license"/>
        </policy>
        <policy context="default">
                <deny own="at.beatthescore.license"/>
        <deny send_destination="at.beatthescore.license"/>
        </policy>
</busconfig>
\end{lstlisting}

Policies have a higher priority the higher it is placed inside the {\tt busconfig} element. The {\tt policy} element contains information about what can be allowed and what has to be denied. Attributes like {\tt user}, {\tt group}, and {\tt context} define which users or groups of users the specific policies apply to. A policy with a {\tt context} value of {\tt "default"} applies to every user.
\\
Since access to the interface for every user is denied the policy configuration makes sure to "punch holes" to allowed users before access gets locked down completely. This is achieved through the two topmost {\tt policy} elements. In this example the {\tt "at.beatthescore.license"} defines the name space on the communication bus. The two topmost policies specify:
\begin{enumerate}
\item{only the user of the backend daemon should be allowed to own the name space for exposing methods on the bus}
\item{only the frontend's user should be allowed to access methods on the bus in the same name space}
\end{enumerate}

Also, by default the D-Bus daemon provides two main buses:
\begin{description}
\item[System bus:] \hfill \\ The system bus allows processes to access methods of other processes which are not running in the session.
\item[Session bus:] \hfill \\ This bus allows processes which are owned by a specific user session to communicate.
\end{description}

\section{Customer login}
Customers create user accounts by logging into the site using Mozilla Persona. Mozilla Persona  \footnote{\label{foot:2}http://persona.org} is an implementation of the BrowserID protocol and allows us to create a login system without having to save login data persistently on the server. Through the use of Persona we are confident that passwords of users won't get compromised in case of a security breach on the server. The email address of the customer is used as the only identifier globally in the database.

The login happens by clicking the respective button on the website.
\\
\includegraphics[width=1.0\textwidth]{bts_login}
\\
When the customer enters his email address he can be authorized to use the website using one of the 3 following ways:

\begin{description}
\item [Authorization using identity provider:] If the email provider has implemented an identity provider the authorization will continue seamlessly. \hfill \\
\item [Authorization using OpenID:] Persona can authorize the user on the website using OpenID in case the provider does not natively implement an identity provider. Of course this only works if the email provider is also an OpenID provider. \hfill \\
\item [Authorization using email:] As a fall back, in case none of the mentioned authorization methods can be used, the customer will get an email sent to his email address. This email would contain a link through which the used can authorize himself manually. \hfill \\
\end{description}

\includegraphics[width=1.0\textwidth]{bts_authorized}

If the authorization is successful the browser will send an HTTP POST request to {\tt beatthescore.at/verify} where we can further validate assertion data with a verification endpoint. In the case of the Beat The Score website we use Mozilla's own verification endpoint (verifier.login.persona.org). Verification is done in the frontend using Mozilla's browserid Python library.
\begin{lstlisting}
@route('/verify', method='POST')
def verify():
    body = request.POST
    assertion = body['assertion']
    print(assertion)
    data = browserid.verify(assertion, ASSERTION_URL)
    licenseManager.InsertUser(data['email'])
    sess = request.environ.get('beaker.session')
    print(sess)
    sess['email'] = data['email']
    return simplejson.dumps({ 'email' : data['email'] })
\end{lstlisting}

\begin{enumerate}
\item In case the verification fails the server will send back an error to the customers web browser and authorization fails.
\item If the verification is done successfully the email address will get sent to the backend process using {\tt licenseManager.InsertUser(data['email'])}. 
\item In case the backend is unable to parse the provided string the email address will not be inserted into the SQLite database and the authorization fails.
\item If the backend process successfully inserts the customers email address into the database a session cookie will be set to indicate a successful authorization to the customers web browser. 
\end{enumerate}

\section{Selling licenses through PayPal}
To further automate the selling of user licenses to customers we decided to use PayPal. PayPal is used by millions of people around the globe and provides several mechanisms for payment automation. All possible ways to integrate PayPal can be found on their developer portal. \footnote{\label{foot:3}https://developer.paypal.com/}
\\
Some of the options which were discussed while designing the payment workflow were:

\begin{description}
\item [PayPal Mobile SDK] Developers can embed views into their applications to allow users to purchase a license. This mechanism can not apply to Beat The Score as there is no SDK for Qt applications. Also this would require releasing a free version of the software package. \hfill \\
\item [PayPal Instant Payment Notification] Using IPN a merchant can receive an HTTP POST request to a URL of choice as soon as a customer completes a payment transaction to the merchants email address. The merchants web server can then send a request to a PayPal server to check if the POST request actually originated from PayPal and whether the received payment transaction is valid or not. The downside is that this approach is asynchronous by design and one of the design goals was to have the whole purchasing process happen in a single browser window. Also, this fits better into use cases where the customer receives the requested goods via email. \hfill \\
\item [PayPal Express Checkout] Express Checkout provides a way to integrate PayPal into the web sites workflow where a successful purchase would immediately result in showing the licenses activation code, all in the same browser window. This activation code could then be entered into the application's window. \hfill \\
\end{description}

\section{License generation}
As soon as the payment transaction is validated and completed, a license file should be generated for the user. This license file is generated in the backend process and inserted into the database afterwards. The license file is generated by using RSA for asymmetrical encryption and SHA256 hashes.

\begin{enumerate}
\item Concatenate the buyers email address and the date of purchase into a string.
\item Generate a SHA256 hash of the concatenated string and represent it in hexadecimal numbers.
\item Concatenate the original string with a newline character and the hexadecimal hash.
\item Encrypt the result with a private RSA key.
\item Apply Base 64 encoding to the result of the encryption.
\end{enumerate}

The Base 64 encoded blob is then inserted into the database. Afterwards, an activation code is generated that is only valid for 5 minutes. Doing this ensures us that the same activation code cannot be used more than once over a long period of time.

\section{Entering the activation code}
As soon as an activation code is generated it can be entered into the application window, together with the customers email address that was used to authorize the user using Mozilla Persona.
\\
\includegraphics[width=1.0\textwidth]{license_enter}
\\
\begin{enumerate}
\item The application will try to validate the activation code and send an HTTP POST request to the frontend server.
\item The frontend server will pass the activation code to the backend process to check if the activation code is valid, if the code is related to an existing license and if the license actually belongs to the entered email address.
\item If all of those checks pass, the backend will read the license from the database and pass it to the frontend.
\item The frontend will in turn send that data back to the application. 
\item The application then saves the license file locally on the file system and validate the license. If the license validation succeeds the user can start using Beat The Score.
\end{enumerate}